\documentclass[12pt,oneside]{article}
\usepackage[T1]{fontenc}
\usepackage{amsmath,amssymb,array,longtable}
\usepackage[english]{babel}
\usepackage[bookmarks=false,pdfauthor={Tomi Setsu},pdftitle={FunC: Syntax Specification}]{hyperref}
\usepackage{fancyhdr}

\setlength{\headheight}{15.2pt}
\pagestyle{fancy}
\fancyhf{}
\fancyfoot[C]{\thepage}
\renewcommand{\headrulewidth}{0.5pt}
\fancyhead[C]{FunC: Syntax Specification}

\title{FunC: Syntax Specification}
\author{Tomi Setsu}

\begin{document}
\maketitle

\begin{abstract}
This document describes the FunC surface syntax currently accepted by the
compiler in this repository. It is derived directly from
\verb|crypto/func/parse-func.cpp|, \verb|crypto/func/keywords.cpp|, and
\verb|crypto/func/builtins.cpp|, with one small supplementary dependency on the
shared lexer implementation for the exact token and comment rules. The target
compiler version is FunC \verb|0.4.6|.

This text is descriptive of the implementation as it exists today. Where a
token is reserved by the lexer but not consumed by the parser, that is stated
explicitly.
\end{abstract}

\section*{Introduction}
FunC is a statically typed source language that compiles into TVM/Fift-style
assembly. This document focuses on its accepted source syntax rather than on the
generated output code.

At the parser level, a FunC source file is a sequence of top-level directives,
global declarations, constant declarations, and function definitions. The core
expression language supports tuples, tensor-valued expressions, explicit type
ascription, method syntax based on identifiers beginning with \verb|.| or
\verb|~|, conditional expressions, and right-associative assignment.

\clearpage
\tableofcontents

\clearpage
\section{Overview}
The current FunC frontend has three relevant layers:

\begin{enumerate}
\item lexical analysis, implemented by the shared lexer;
\item keyword registration, implemented in \verb|keywords.cpp|;
\item parsing and type-directed expression construction, implemented in
\verb|parse-func.cpp|.
\end{enumerate}

The compiler also injects a small built-in prelude before parsing any user
source. This prelude is defined in \verb|builtins.cpp| and contains the
primitive arithmetic operators, comparison operators, a few constants, basic
slice and builder helpers, and several debugging and dynamic-dispatch helpers.

\section{Lexical Structure}
\subsection{Whitespace and Comments}
Whitespace separates tokens in the usual way. The current lexer recognizes:

\begin{itemize}
\item end-of-line comments beginning with \verb|;;|;
\item nested block comments delimited by \verb|{-| and \verb|-}|.
\end{itemize}

Strings are delimited either by ordinary double quotes or by triple double
quotes. Triple-quoted strings may span multiple lines.

\subsection{Token Boundaries}
The FunC parser instantiates the lexer with \verb|";,()[] ~."| as the active
character specification. In practice, this means:

\begin{itemize}
\item \verb|;|, \verb|,|, \verb|(|, \verb|)|, \verb|[|, and \verb|]| are always
standalone tokens;
\item \verb|.| and \verb|~| act as left-active prefix characters so that tokens
such as \verb|.store_uint| and \verb|~load_uint| are read as single identifiers;
\item most other operators rely on ordinary token separation by whitespace or by
the surrounding punctuation.
\end{itemize}

For this reason, source code normally writes infix operators with whitespace and
writes block braces with ordinary separation, for example:

\begin{verbatim}
if (x) { return (); }
int y = x + 1;
\end{verbatim}

\subsection{Identifiers}
The current lexer does not enforce a narrow identifier regular expression.
Instead, any token that is not recognized as a number or registered keyword
becomes an identifier.

As a result, valid identifiers in practice include names such as:

\begin{verbatim}
recv_internal
op::transfer
slice_empty?
dict_get?
.store_uint
~load_uint
\end{verbatim}

The parser distinguishes three identifier classes:

\begin{itemize}
\item ordinary identifiers;
\item \emph{dot identifiers}, whose first character is \verb|.|;
\item \emph{tilde identifiers}, whose first character is \verb|~|.
\end{itemize}

Dot and tilde identifiers are used by method-call syntax; see
Section~\ref{sec:method-calls}.

Back-quoted tokens are also supported by the shared lexer. A token of the form
\verb|`...`| is first read as a raw token and then classified as if the quotes
were not present. This is mainly useful when a name would otherwise be hard to
tokenize.

\subsection{Integer Literals}
An integer literal is either:

\begin{itemize}
\item a decimal literal, such as \verb|0|, \verb|17|, or \verb|-3|;
\item a hexadecimal literal beginning with \verb|0x|, such as \verb|0xff|.
\end{itemize}

The parser converts integer literals into signed 257-bit values and rejects
anything that does not fit.

\subsection{String Literals}
The parser accepts ordinary quoted strings and triple-quoted strings, each with
an optional one-letter suffix. The suffix determines the semantic class of the
literal:

\begin{center}
\begin{tabular}{|>{\ttfamily}c|m{0.67\textwidth}|}
\hline
suffix & meaning \\ \hline
none & slice constant built from the raw bytes of the literal \\ \hline
s & slice constant; the literal body is parsed as a hexadecimal bitstring \\ \hline
a & slice constant; the literal body is parsed as a standard address \\ \hline
u & integer constant obtained from the raw bytes of the literal \\ \hline
h & integer constant equal to the first 32 bits of SHA-256 of the literal \\ \hline
H & integer constant equal to the full 256-bit SHA-256 of the literal \\ \hline
c & integer constant equal to CRC32 of the literal \\ \hline
\end{tabular}
\end{center}

No other string suffix is recognized by the current parser.

\subsection{Reserved Keywords and Operators}
The current keyword table reserves the following words:

\begin{verbatim}
return  var  repeat  do  while  until  try  catch
if  ifnot  then  else  elseif  elseifnot
int  cell  slice  builder  cont  tuple  type  forall
extern  global  asm  impure  inline  inline_ref  auto_apply
method_id  operator  infix  infixl  infixr  const
#pragma  #include
\end{verbatim}

The following multi-character operators are also reserved:

\begin{verbatim}
==  !=  <=  >=  <=>
<<  >>  ~>>  ^>>
~/  ^/  ~%  ^%  /%
+=  -=  *=  /=  ~/=  ^/=  %=  ~%=  ^%=
<<=  >>=  ~>>=  ^>>=
&=  |=  ^=
\end{verbatim}

In addition, the single characters below are registered as keywords in their
own right:

\begin{verbatim}
+  -  *  /  %  ?  :  ,  ;  (  )  [  ]  {  }  =  _
<  >  &  |  ^  ~
\end{verbatim}

\section{Types}
\label{sec:types}
\subsection{Atomic Types}
The parser recognizes the following atomic types:

\begin{verbatim}
int  cell  slice  builder  cont  tuple  type
\end{verbatim}

In addition, \verb|_| and \verb|var| both denote a fresh type hole in type
position.

\subsection{Type Constructors}
The parser uses the following grammar for type expressions:

\begin{verbatim}
TypeExpr  ::= AtomicType
            | "_"
            | "var"
            | "()"
            | "(" TypeExpr ")"
            | "(" TypeExpr "," TypeExpr { "," TypeExpr } ")"
            | "[]"
            | "[" TypeExpr { "," TypeExpr } "]"
            | TypeExpr "->" TypeExpr
\end{verbatim}

The meaning of the constructors is:

\begin{itemize}
\item \verb|()| is the unit tensor type;
\item \verb|(T1, T2, ..., Tn)| is a tensor type of width \(n\);
\item \verb|[T1, T2, ..., Tn]| is a tuple type, still represented as a single
runtime value;
\item \verb|A -> B| is a function type, parsed right-associatively.
\end{itemize}

The distinction between tensors and tuples is important. A tensor type denotes
multiple stack values, while a tuple type denotes one TVM tuple value.

Examples:

\begin{verbatim}
int
(int, slice)
[int, cell]
int -> int
(int, slice) -> [int, slice]
\end{verbatim}

\subsection{Universally Quantified Type Variables}
Source-level universal quantification appears only as a prefix of a function
definition:

\begin{verbatim}
forall X, Y -> [X, Y] pair(X x, Y y) asm "PAIR";
\end{verbatim}

Each bound identifier becomes a type variable visible in the function
signature. The optional keyword \verb|type| may appear before each such
identifier, but it is not required:

\begin{verbatim}
forall type X, type Y -> ...
forall X, Y -> ...
\end{verbatim}

\section{Top-Level Source Structure}
\label{sec:toplevel}
A source file is parsed as:

\begin{verbatim}
SourceItem ::= Pragma
             | Include
             | GlobalDecls
             | ConstDecls
             | FunctionDef

Source     ::= { SourceItem }
\end{verbatim}

\subsection{Pragmas}
The directive syntax is:

\begin{verbatim}
Pragma ::= "#pragma" identifier ... ";"
\end{verbatim}

The current parser recognizes these pragma names:

\begin{itemize}
\item \verb|version|;
\item \verb|not-version|;
\item \verb|test-version-set|;
\item \verb|allow-post-modification|;
\item \verb|compute-asm-ltr|.
\end{itemize}

\verb|version| and \verb|not-version| accept semver-style constraints with the
comparators \verb|=|, \verb|>|, \verb|>=|, \verb|<|, \verb|<=|, and \verb|^|.
\verb|test-version-set| replaces the compiler version string for test purposes.

\subsection{Includes}
The include syntax is:

\begin{verbatim}
Include ::= "#include" string-literal ";"
\end{verbatim}

Included paths are resolved relative to the current file before being passed to
the compiler's file-reading callback. The implementation suppresses repeated
inclusion of the same real path.

\subsection{Global Variables}
The syntax for global variable declarations is:

\begin{verbatim}
GlobalDecls ::= "global" GlobalDecl { "," GlobalDecl } ";"
GlobalDecl  ::= [ TypeExpr ] identifier
\end{verbatim}

If the type is omitted, or if the declaration begins with \verb|_|, the parser
creates a fresh type hole and lets later usage constrain it.

Examples:

\begin{verbatim}
global int seqno;
global cell state, slice config;
global _ scratch;
\end{verbatim}

\subsection{Constants}
The syntax for constant declarations is:

\begin{verbatim}
ConstDecls ::= "const" ConstDecl { "," ConstDecl } ";"
ConstDecl  ::= [ "int" | "slice" ] identifier "=" ConstExpr
\end{verbatim}

The current implementation accepts:

\begin{itemize}
\item integer literals;
\item slice/string literals;
\item compile-time integer expressions that can be reduced by the compiler into
a single constant instruction.
\end{itemize}

Constants are therefore much more restrictive than ordinary expressions.

Examples:

\begin{verbatim}
const int one = 1;
const slice prefix = "deadbeef"s;
const int answer = 6 * 7;
\end{verbatim}

\section{Function Definitions}
\label{sec:function-defs}
\subsection{Grammar}
The surface grammar for function definitions is:

\begin{verbatim}
FunctionDef ::= [ ForallClause ]
                TypeExpr identifier FormalArgs
                [ "impure" ]
                [ "inline" | "inline_ref" ]
                [ MethodIdClause ]
                ( ";" | Block | AsmBody )

ForallClause ::= "forall" [ "type" ] identifier
                 { "," [ "type" ] identifier } "->"

FormalArgs   ::= "(" [ FormalArg { "," FormalArg } ] ")"
FormalArg    ::= "_"
               | [ TypeExpr ] [ identifier ]

MethodIdClause ::= "method_id"
                 | "method_id" "(" string-literal ")"
                 | "method_id" "(" integer-literal ")"
\end{verbatim}

The order of the optional modifiers is significant because it is hard-coded in
the parser:

\begin{enumerate}
\item \verb|impure|;
\item one of \verb|inline| or \verb|inline_ref|;
\item \verb|method_id|.
\end{enumerate}

\subsection{Forward Declarations}
If a function definition ends with \verb|;|, it is a declaration without a body:

\begin{verbatim}
int helper(int x);
\end{verbatim}

Repeated declarations must unify with the already known function type.

\subsection{Ordinary Bodies}
An ordinary function body is a block:

\begin{verbatim}
int add(int x, int y) {
  return x + y;
}
\end{verbatim}

If control can fall off the end of the function body, the parser inserts an
implicit \verb|return ();| and requires the declared return type to unify with
unit.

\subsection{Assembler Bodies}
\label{sec:asm-bodies}
An assembler-bodied function uses the syntax:

\begin{verbatim}
AsmBody ::= "asm"
            [ "(" [ ArgOrder ] [ "->" RetOrder ] ")" ]
            string-literal { string-literal }
            ";"
\end{verbatim}

The optional permutation clause has two roles:

\begin{itemize}
\item \verb|ArgOrder| lists formal parameter names in the order expected by the
assembly implementation;
\item \verb|RetOrder| lists integer indices \verb|0 .. width-1| describing how
multiple returned stack values are reordered.
\end{itemize}

Example:

\begin{verbatim}
forall X -> (X, tuple) list_next(tuple list) asm( -> 1 0) "UNCONS";
\end{verbatim}

The parser imposes two important static restrictions on \verb|asm| functions:

\begin{itemize}
\item at most 16 formal arguments;
\item a fixed-width return type and fixed-width argument types.
\end{itemize}

\subsection{Method Identifiers}
If \verb|method_id| is present without an explicit argument, the compiler uses
the function name. If its argument is a string, the compiler computes
\texttt{crc16(name) | 0x10000}. If its argument is an integer, that integer is
used directly.

Examples:

\begin{verbatim}
int seqno() method_id { ... }
int get_public_key() method_id("pubkey") { ... }
int my_getter() method_id(123456) { ... }
\end{verbatim}

\section{Statements}
\label{sec:statements}
\subsection{Grammar}
Function bodies are made of the following statements:

\begin{verbatim}
Stmt ::= "return" Expr ";"
       | Block
       | ";"
       | "repeat" Expr Block
       | "while" Expr Block
       | "do" Block "until" Expr
       | "try" Block "catch" Expr Block
       | IfStmt
       | Expr ";"

Block ::= "{" { Stmt } "}"

IfStmt ::= ("if" | "ifnot") Expr Block
           [ "else" Block
           | ("elseif" | "elseifnot") Expr Block [ ... ] ]
\end{verbatim}

\subsection{Return Statements}
The current parser does \emph{not} support a bare \verb|return;|. Unit return
must be written explicitly:

\begin{verbatim}
return ();
\end{verbatim}

\subsection{Condition and Loop Types}
The parser requires the following expressions to unify with \verb|int| and to
produce a single stack value:

\begin{itemize}
\item \verb|if| and \verb|ifnot| conditions;
\item \verb|elseif| and \verb|elseifnot| conditions;
\item \verb|while| conditions;
\item \verb|do ... until| conditions;
\item \verb|repeat| iteration counts.
\end{itemize}

\subsection{Try/Catch}
A \verb|catch| clause binds an expression pattern whose type must unify with a
two-component tensor \verb|(X, int)|. In normal source code this is usually
written as two variables or holes:

\begin{verbatim}
try {
  ...
} catch (exc_arg, exc_code) {
  ...
}
\end{verbatim}

\section{Expressions}
\label{sec:expressions}
\subsection{Primary Forms}
The highest-precedence expression forms are:

\begin{verbatim}
Primary ::= "(" ")"
          | "(" Expr ")"
          | "(" Expr "," Expr { "," Expr } ")"
          | "[" "]"
          | "[" Expr { "," Expr } "]"
          | integer-literal
          | string-literal
          | "_"
          | "var"
          | "int" | "cell" | "slice" | "builder" | "cont" | "tuple" | "type"
          | identifier
\end{verbatim}

Their meaning is:

\begin{itemize}
\item \verb|(e1, e2, ..., en)| constructs a tensor-valued expression;
\item \verb|[e1, e2, ..., en]| constructs a TVM tuple value;
\item \verb|_| constructs a hole pattern suitable for destructuring on the left
of assignment;
\item a type keyword or type identifier in expression position introduces an
explicit type application, described below.
\end{itemize}

\subsection{Application and Explicit Type Application}
After a primary expression, the parser repeatedly accepts another primary
expression. The resulting form means one of two things:

\begin{enumerate}
\item if the left expression is a type expression, the form is an explicit type
application used for typed patterns and declarations;
\item otherwise, it is ordinary function application.
\end{enumerate}

Examples:

\begin{verbatim}
f(x)
f(x, y)
int x = 1;
var msg_seqno = cs~load_uint(32);
(int, slice) (n, s) = (1, "ab"s);
\end{verbatim}

This rule is the reason local variable introduction is syntax-driven. A fresh
local variable is created by writing it under an explicit type pattern such as
\verb|int x|, \verb|var x|, or a tuple/tensor pattern containing such names.
A bare unknown name on the left of \verb|=| is \emph{not} the declaration form
implemented by the current parser.

\subsection{Method Calls}
\label{sec:method-calls}
After an expression, the parser accepts any number of special identifiers whose
text begins with \verb|.| or \verb|~|. The general form is:

\begin{verbatim}
receiver.special_identifier argument
\end{verbatim}

In source code the special identifier is normally attached directly to the
receiver:

\begin{verbatim}
cs~load_uint(32)
begin_cell().store_uint(x, 32).end_cell()
\end{verbatim}

The difference between the two prefixes is:

\begin{itemize}
\item \verb|.name| requires the receiver to be an rvalue;
\item \verb|~name| requires the receiver to be an lvalue and is compiled through
a ``modify-and-return'' pattern.
\end{itemize}

If a \verb|~name| symbol is not defined but the corresponding plain
\verb|name| symbol exists, the parser falls back to the plain symbol.

\subsection{Operator Precedence}
The parser implements the following precedence levels, from highest to lowest:

\begin{center}
\begin{tabular}{|l|m{0.63\textwidth}|}
\hline
level & forms \\ \hline
100 & primary expressions \\ \hline
90 & repeated application and explicit type application \\ \hline
80 & repeated method syntax using identifiers beginning with dot or tilde \\ \hline
75 & prefix tilde \\ \hline
30 & multiplicative operators, rounded division/modulus operators, and bitwise and \\ \hline
20 & prefix minus, then additive operators, bitwise or, and bitwise xor \\ \hline
17 & shift operators \\ \hline
15 & one comparison operator \\ \hline
13 & conditional operator \\ \hline
10 & assignment and compound assignment operators \\ \hline
\end{tabular}
\end{center}

Important implementation details:

\begin{itemize}
\item comparison operators are not chainable in the current parser;
\item assignment is right-associative;
\item prefix \verb|-| is parsed at level 20, so its operand is a level-30
expression rather than a level-75 expression;
\item negative numeric literals such as \verb|-1| are usually lexed as single
number tokens, which avoids ambiguity for the common literal case.
\end{itemize}

\subsection{Assignments and Compound Assignments}
The assignment grammar is:

\begin{verbatim}
AssignExpr ::= LValue "=" Expr
             | LValue "+=" Expr
             | LValue "-=" Expr
             | LValue "*=" Expr
             | LValue "/=" Expr
             | LValue "~/=" Expr
             | LValue "^/=" Expr
             | LValue "%=" Expr
             | LValue "~%=" Expr
             | LValue "^%=" Expr
             | LValue "<<=" Expr
             | LValue ">>=" Expr
             | LValue "~>>=" Expr
             | LValue "^>>=" Expr
             | LValue "&=" Expr
             | LValue "|=" Expr
             | LValue "^=" Expr
\end{verbatim}

Assignments are expressions, not separate statements. They therefore may appear
wherever an rvalue expression is allowed, although they are most often used as
statement bodies.

\section{Compiler Prelude and Built-ins}
\label{sec:builtins}
Before parsing user code, the compiler defines a small built-in prelude. The
surface language uses many of these names indirectly via operator syntax or
method syntax.

\subsection{Surface Operators Backed by Built-in Symbols}
The following built-in function names are the parser's canonical targets for the
surface operators:

\begin{verbatim}
_+_   _-_   -_   _*_   _/_   _~/_   _^/_   _%_   _~%_   _^%_   _/%_
_<<_  _>>_  _~>>_  _^>>_  _&_  _|_  _^_  ~_
_==_  _!=_  _<_  _>_  _<=_  _>=_  _<=>_
\end{verbatim}

Compound assignments use a parallel family of internal names:

\begin{verbatim}
^_+=_   ^_-=_   ^_*=_   ^_/=_   ^_~/=_   ^_^/=_
^_%=_   ^_~%=_  ^_^%=_  ^_<<=_  ^_>>=_   ^_~>>=_  ^_^>>=_
^_&=_   ^_|=_   ^_^=_
\end{verbatim}

\subsection{Named Built-ins}
The current named built-ins are grouped as follows.

Arithmetic helpers:
\begin{verbatim}
divmod  ~divmod  moddiv  ~moddiv
muldiv  muldivr  muldivc  muldivmod
\end{verbatim}

Constants and null checks:
\begin{verbatim}
true  false  nil  Nil  null?
\end{verbatim}

Exceptions:
\begin{verbatim}
throw  throw_if  throw_unless
throw_arg  throw_arg_if  throw_arg_unless
\end{verbatim}

Slice and builder primitives:
\begin{verbatim}
load_int  load_uint  preload_int  preload_uint
store_int  store_uint  ~store_int  ~store_uint
load_bits  preload_bits
\end{verbatim}

Tuple access:
\begin{verbatim}
int_at  cell_at  slice_at  tuple_at  at
\end{verbatim}

No-op and debugging helpers:
\begin{verbatim}
touch  ~touch  touch2  ~touch2  ~dump  ~strdump
\end{verbatim}

Dynamic calls:
\begin{verbatim}
run_method0  run_method1  run_method2  run_method3
\end{verbatim}

\subsection{Notes}
The nearby source code contains a commented-out definition for \verb|null|, but
\verb|null| itself is \emph{not} currently injected by \verb|builtins.cpp|.
Only \verb|null?| is part of the compiler prelude.

\section{Reserved but Currently Unused Keywords}
The following tokens are reserved by \verb|keywords.cpp| but are not consumed by
the current FunC parser:

\begin{verbatim}
then
extern
auto_apply
operator
infix
infixl
infixr
\end{verbatim}

They therefore cannot be used as ordinary identifiers, but they also have no
surface syntax today.

\appendix
\section{Summary Grammar}
This appendix collects the main grammar fragments in one place.

\begin{verbatim}
Source     ::= { SourceItem }
SourceItem ::= Pragma | Include | GlobalDecls | ConstDecls | FunctionDef

Pragma     ::= "#pragma" identifier ... ";"
Include    ::= "#include" string-literal ";"

GlobalDecls ::= "global" GlobalDecl { "," GlobalDecl } ";"
GlobalDecl  ::= [ TypeExpr ] identifier

ConstDecls ::= "const" ConstDecl { "," ConstDecl } ";"
ConstDecl  ::= [ "int" | "slice" ] identifier "=" ConstExpr

FunctionDef ::= [ ForallClause ]
                TypeExpr identifier FormalArgs
                [ "impure" ]
                [ "inline" | "inline_ref" ]
                [ MethodIdClause ]
                ( ";" | Block | AsmBody )

ForallClause ::= "forall" [ "type" ] identifier
                 { "," [ "type" ] identifier } "->"

FormalArgs ::= "(" [ FormalArg { "," FormalArg } ] ")"
FormalArg  ::= "_"
             | [ TypeExpr ] [ identifier ]

MethodIdClause ::= "method_id"
                 | "method_id" "(" string-literal ")"
                 | "method_id" "(" integer-literal ")"

AsmBody ::= "asm"
            [ "(" [ ArgOrder ] [ "->" RetOrder ] ")" ]
            string-literal { string-literal }
            ";"

Block ::= "{" { Stmt } "}"

Stmt ::= "return" Expr ";"
       | Block
       | ";"
       | "repeat" Expr Block
       | "while" Expr Block
       | "do" Block "until" Expr
       | "try" Block "catch" Expr Block
       | IfStmt
       | Expr ";"

IfStmt ::= ("if" | "ifnot") Expr Block
           [ "else" Block
           | ("elseif" | "elseifnot") Expr Block [ ... ] ]

TypeExpr ::= AtomicType
           | "_"
           | "var"
           | "()"
           | "(" TypeExpr ")"
           | "(" TypeExpr "," TypeExpr { "," TypeExpr } ")"
           | "[]"
           | "[" TypeExpr { "," TypeExpr } "]"
           | TypeExpr "->" TypeExpr
\end{verbatim}

\section{Expression Precedence Summary}
\begin{verbatim}
100  primary expressions
 90  repeated application / explicit type application
 80  repeated .method and ~method application
 75  prefix ~
 30  *  /  ~/  ^/  %  ~%  ^%  /%  &
 20  prefix - ; then +  -  |  ^
 17  <<  >>  ~>>  ^>>
 15  ==  !=  <  >  <=  >=  <=>
 13  ?:
 10  =  +=  -=  *=  /=  ~/=  ^/=  %=  ~%=  ^%=  <<=  >>=  ~>>=  ^>>=  &=  |=  ^=
\end{verbatim}

\end{document}
